%% kapa: exam=2 round=21 year=2010 period=1 subject=appraisal_practice subject_ko=감정평가실무 seq=2 number=2 points=20
%% source: Q-Net 감정평가사 기출문제 (cst00309) artlSeq=5024183 file=extracted/1008/5024183/2010년 제21회 감정평가사 2차시험 기출문제/1교시.hwp
%% conversion: kordoc 4.24.51 -> tools/gen2.py (auto); review=proofread-sonnet; flags=none
\begin{longquestion}{【문제 2】}
이대한씨는 甲구 乙동에 근린생활시설을 소유하고 있다. 이 건물의 1층은 편의점으로, 2층은 주거용으로 사용되고 있다. 대상 토지는 제1종 주거지역내 정방형의 대(지목)로서 150m\textsuperscript{2}중 50m\textsuperscript{2}가 공익사업인 자동차전용 도로개설사업에 편입되어 보상협의요청서를 받았으나 토지에 대하여는 협의에 응하지 않고 재결의 신청을 청구하였는바, 사업시행자로 부터 재결서 정본을 수령하였다. 재결에 불복한 이대한씨는 이의신청과 동시에 잔여지의 손실에 대하여도 손실보상을 청구하기로 하였다. 다음의 자료를 검토, 분석 후 물음에 답하시오. (20점) \par
물음(1) 잔여지(100m\textsuperscript{2})와 관련하여 다음 질문에 답하시오. (10점) \par
1) 잔여지에 대한 손실보상을 받기 위한 요건을 약술하시오. \par
2) 잔여지 손실보상의 종류와 각각의 종류에 따른 보상액 산정방법을 설명하되, 자료를 활용하여 산정 가능한 범위 내에서 이대한씨가 청구할 수 있는 손실보상액을 산정하시오. \par
물음(2) 공정감정평가법인 소속 김한국 감정평가사가 이대한씨의 영업손실보상액을 산정(협의보상)하기 위하여 수집하였을 것으로 판단되는 제반자료 및 조사사항을 “공익사업을 위한 토지 등의 취득 및 보상에 관한 법률” 및 동법 시행규칙, “영업손실 보상평가지침”등과 관련하여 설명하시오. (10점) \par
<자료 1> 손실보상금 지급내역서 (요약) \par
{\small\centering
\begin{tabular}{|c|c|c|c|c|c|}
\hline
\makecell{소재지} & \makecell{지번} & \makecell{면적} & \makecell{편입면적} & \makecell{단가(원/m\textsuperscript{2})} & \makecell{금액(원)} \\
\hline
\makecell{甲구 乙동} & \makecell{123-1} & \makecell{150} & \makecell{50} & \makecell{1,500,000} & \makecell{75,000,000} \\
\hline
\end{tabular}\par}

※ 건물등 지장물과 영업손실의 보상금은 협의를 위한 보상협의요청시 수령하였음. \par
<자료 2> 대상 부동산의 상황 등 \par
1) 잔여지 100m\textsuperscript{2}는 자동차 전용도로가 잔여지의 전면부를 통과하여 맹지가 되어 건축이 불가능한 상태임. \par
2) 가격조사 결과 공익사업시행 후, 잔여지 매매가능가액은 700,000원/m\textsuperscript{2} 내외로 파악되었음. \par
3) 도로개설공사 완료일 : 2010년 9월 5일 \par
4) 본 건의 사업인정고시일 : 2009년 1월 1일 \par
5) 본 사업으로 인한 인근지역의 지가변동은 없는 것으로 조사되었음. \par
\end{longquestion}
